Model Context Protocol (MCP) is emerging as a ``REST API for AI agents,'' standardizing how models call tools and consume structured results. However, MCP’s control plane is JSON-RPC, and JSON is structurally inefficient for large tabular tool outputs: it inflates payload bytes, stresses context budgets, and blocks incremental consumption. We present \system, a control/data-plane split for MCP tool results. Tools return a small MCP response (a \texttt{CallToolResult} envelope) containing human-readable text plus a machine-readable descriptor that points to a data plane served over streamable HTTP. For large tables, the data plane delivers Parquet either as a whole blob or as a length-prefixed chunk stream, enabling time-to-first-rows, early termination, and bounded client memory. \system includes a format-selection framework that uses optimization targets (min bytes, min latency, or min time-to-first-rows), server-provided size hints, and optional online bandit learning; within the Parquet arm it can also choose codecs/encodings.\ We evaluate on (i) controlled synthetic grids and TPC-DS \texttt{catalog\_sales} slices to characterize size scaling and decode proxies, and (ii) BIRD dev SQLite query results to measure real workload distributions. On synthetic grids, JSON transfers are 2.7--2.9$\times$ larger than Parquet blobs. On BIRD full dev (1534 queries), hint-aware selection reduces aggregate wire volume from 43.4\,MB (JSON baseline) to 6.4\,MB (6.8$\times$), while preserving millisecond-scale localhost fetch latency; the additional ``describe'' round-trip can be avoided with a one-shot server-side selector. These results suggest that separating MCP control from a tabular data plane is a practical path to scalable agent-first data access.